\documentclass[../main.tex]{subfiles}
\begin{document}

El ransomware se ha convertido en una de las amenazas más graves para
organizaciones de todo tipo. Los ataques modernos no se limitan a cifrar
los datos de la víctima: incluyen sistemáticamente la destrucción de las
copias de seguridad antes del cifrado, eliminando la principal vía de
recuperación. Frente a este escenario, la industria propone el
almacenamiento inmutable como mecanismo de protección, pero la literatura
académica carece de comparativas experimentales que cuantifiquen su impacto
real sobre la recuperación.

Este Trabajo de Fin de Grado presenta el diseño, implementación y evaluación
de un laboratorio de simulación de ransomware basado íntegramente en
contenedores Docker. El entorno reproduce una infraestructura empresarial
mínima compuesta por un servidor de ficheros, una \textit{workstation}, un
servidor de \textit{backups} SFTP, un sistema de almacenamiento de objetos
con Object Lock (MinIO) y un SIEM completo (Wazuh~4.7.5 con OpenSearch).

El diseño experimental sigue un esquema factorial de cuatro escenarios de
ataque por tres estrategias de recuperación y veinte repeticiones,
totalizando 240~ejecuciones completamente automatizadas y ejecutadas en orden
aleatorio. Los escenarios cubren las tácticas más relevantes de la matriz
MITRE ATT\&CK: cifrado masivo de ficheros (T1486), destrucción activa de
copias de seguridad (T1490), movimiento lateral por SSH (T1021.004) y
exfiltración de datos (T1041). Las tres estrategias evaluadas son: ausencia
de copia de seguridad, repositorio Restic sobre SFTP convencional y
repositorio Restic sobre MinIO con Object Lock COMPLIANCE.

El hallazgo central es de disponibilidad: ante la destrucción activa del
repositorio, la estrategia SFTP resulta irrecuperable en el 100\,\% de los
casos, mientras que la combinación de Object Lock y una política de mínimo
privilegio mantiene una tasa de éxito del 100\,\% (test exacto de Fisher
$p < 0.001$). Además, cuando se mide de forma aislada la velocidad del
\textit{backend}, la restauración desde el almacenamiento inmutable resulta
un 39\,\% más rápida que desde SFTP, con un tamaño del efecto máximo
($\delta$ de Cliff $=1.0$). Los tiempos de detección mediante Wazuh, medidos
a partir del instante real de la alerta, oscilan entre 0.4 y 11.2~segundos,
muy por debajo del umbral de un minuto y en marcado contraste con los tiempos
de permanencia de meses reportados por la industria para incidentes no
instrumentados.

El laboratorio, con detección por monitorización de integridad de ficheros (FIM) y eventos estructurados, reglas de
correlación Wazuh \textit{custom} y recuperación automatizada verificada por
\textit{hash}, se publica como software libre para facilitar la reproducción
independiente de los experimentos.

\end{document}
